\documentclass[11pt]{article}
\usepackage[utf8]{inputenc}	% Para caracteres en español
\usepackage{amsmath,amsthm,amsfonts,amssymb,amscd}
\usepackage{multirow,booktabs}
\usepackage[table]{xcolor}
\usepackage{fullpage}
\usepackage{lastpage}
\usepackage{enumitem}
\usepackage{fancyhdr}
\usepackage{mathrsfs}
\usepackage{wrapfig}
\usepackage{setspace}
\usepackage{calc}
\usepackage{multicol}
\usepackage{cancel}
\usepackage[retainorgcmds]{IEEEtrantools}
\usepackage[margin=3cm]{geometry}
\usepackage{amsmath}
\newlength{\tabcont}
\setlength{\parindent}{0.0in}
\setlength{\parskip}{0.05in}
\usepackage{empheq}
\usepackage{framed}
\usepackage[most]{tcolorbox}
\usepackage{xcolor}
\colorlet{shadecolor}{orange!15}
\parindent 0in
\parskip 12pt
\geometry{margin=1in, headsep=0.25in}
\theoremstyle{definition}
\newtheorem{defn}{Definition}
\newtheorem{reg}{Rule}
\newtheorem{exer}{Exercise}
\newtheorem{note}{Note}

% Each lecture starts on a new page
\usepackage{titlesec}
\newcommand{\sectionbreak}{\clearpage}


\begin{document}
\setcounter{section}{0}
\title{Econ 457 Lecture Notes}

\thispagestyle{empty}

% Some tools:
% Equation: \begin{equation}
% Notes: \begin{note}
% Shaded area: \begin{shaded}
% Definition: \begin{defn}



\begin{center}
{\LARGE \bf Econ 457}\\
{\bf Lecture Notes}\\
\end{center}

\section{Lecture 1}
\subsection{GDP}

\begin{itemize}
\item \textbf{Gross Domestic Product (GDP)} is the sum of value added in the economy during a given period.   It includes consumption (goods and services), investment (structures, equipment, and inventories), and government expenditures.  GDP does not \textit{directly} include the value of financial assets.   There is a small, indirect contribution through the fees charged by asset managers (a payment for their service), which are usually based on the value of financial assets.
\end{itemize}


\section{Lecture 2: Ch 5, 1: Measurement}


    \begin{shaded}
    \subsection{Practice: EAR, APR, $r_{cc}$}
    \begin{enumerate}
        \item You purchase a security on January 1st for \$98.  The security produces quarterly income of \$2.  You sell the security after 9 months for \$105.   Calculate the following:
    \begin{itemize}
        \item HPR: ($2\times3 + (105-98)) / 98 = 13\%$
        \item EAR: $(1+13\%)^{\frac{1}{0.75}} - 1 = 17\%$
    \end{itemize}
    \item An instrument has a APR of 15\% and compounds monthly.  Calculate the following:
        \begin{itemize}
        \item Monthly compounding rate (in \%): $15/12 = 1.25\%$
        \item EAR: $(1+1.25\%)^{12} -1 = 16\%$
        \item $r_{cc}: e^{r_{cc}} = 1.16 \\ r_{cc} = ln(1.16) = 14.8\%$
    \end{itemize}
    \item An investment has an EAR of 10\% and compounds monthly. Calculate the following:
    \begin{itemize}
        \item APR: $10/12 = 0.88\%$
        \item $r_{cc}:$ 
    \end{itemize}
    \end{enumerate}
    \end{shaded}

    \begin{shaded}
    \subsection{Practice: Geomtric v Arithmetic Mean}
    \begin{enumerate}
    \item Stock ABC doesn't pay dividends.   In year 1 the price of ABC declines by 20\%.   
    In year 2 the price rebounds by 20\%. What is the arithmetic average of the annual returns?   What is the geometric average?      Which measure is more appropriate?\\
    \vspace{1em}
    
   \textit{ Arithmetic average = $0\%$
    Geometric average = $\sqrt{(1+0.2) \times (1-0.2)} = \sqrt{1.2 \times 0.8} = \sqrt{0.96} = 0.98 - 1 = -2\%$}
    \vspace{1em}

    \item   You invest \$1 million at the beginning of 2028 in a stock-index fund. 
    If the rate of return in 2028 is -40\%, what rate of return in 2029 will be necessary for your 
    portfolio to recover to its original value.\\
    \vspace{1em}
    
    \textit{Answer: $\$400,000 / \$600,000 = 1.66 \implies 66\%$ return}
    \end{enumerate}

    \end{shaded}

    \subsection{Logs}
    In finance the analysis is often of the natural log of the price.   Here are a few considerations:
    \begin{itemize}
    \item The difference in the log price series is equal to the log of the gross return series
    \begin{align*}
    ln(P_t) - ln(P_{t-1}) &= ln\left(\frac{P_t}{P_{t-1}}\right)\\
    &= ln(1+r_t)
    \end{align*}
    \item In practice, the difference of the logs is very close to the returns.   
    $$ln(1+r_t) \approx r_t$$
    This is due to the Taylor Rule.   The Taylor Expansion of $f(x)$ evaluated at $x=a$ is
    \begin{align*}
        f(x) &= \sum_{n=0}^{\infty}\frac{f^{(n)}(a)}{n!}(x-a)^2 \\
            &= f(a) + \frac{f'(a)}{1!}(x-a) + \frac{f''(a)}{2!}(x-a)^2 + \frac{f'''(a)}{3!}(x-a)^3 + \dots \\
    \end{align*}
    For the special case that $a=0$
        $$f(x) = f(0) + \frac{f'(0)}{1!}x + \frac{f''(0)}{2!}x^2 + \frac{f'''(0)}{3!}x^3 + \dots$$
    When $f(x) = ln(1+x)$, then $f'(x) = \frac{1}{1+x}$ and $f''(x) = \frac{-1}{(1+x)^2}$, etc
    \begin{align*}
        ln(1+x) &= ln(1) + x -\frac{x^2}{2} + \frac{x^3}{3} - \frac{x^4}{4} + \dots\\
         &= x
    \end{align*}
    \item One reason to prefer the log of the price is that the logs are added, which is needed for the normal distribution results
    \item A second reason to prefer the log of the price is that the distribution of future (log) of price will also be normally distributed.
    \item A final reason to prefer the log of the price is that this linearizes graphs, which makes it easier to see trends.
    \begin{itemize}
    \item If the y-axis is in a log scale, then an increment of 1 represents an increase of 171\%.  This follows because if $ln(P_1) - ln(P_0) = 1$ then $ln(\frac{P_1}{P_0}) = 1$ then $ \frac{P_1}{P_0} = e = 2.71$ then $P_1 = 2.71P_0$ and $P_1 - P_0  = 1.71P_0$.   Note that an increase of 171\% is the same as a continuously compounded return of 100\%.
    \item The slope of the log-price graph is equal to the average continuously compounded return.
    \end{itemize}
    \end{itemize}

    \begin{shaded}
    \subsection{Practice: Log Scale}
    \begin{itemize}
        \item If the natural log of a total return series increase by 1 point over a 1-year period, what is the HPR for that period?   What is the EAR?\\
            \vspace{1em}
            $$ln{\frac{P_2}{P_0}} = 1$$
            Exponential of both sides
            $$\frac{P_2}{P_0} = e$$
            $$HPR = e - 1 = 171\%$$
            $$EAR = HPR$$
            \item If the natural log of a total return series increase by 1 point over a 10-year period, what is the HPR for that period?   What is the EAR?  What is the $r_{cc}$\\
            \vspace{1em}
            $$ln{\frac{P_{10}}{P_0}} = 1$$
            Exponential of both sides
            $$\frac{P_{10}}{P_0} = e$$
            $$HPR = e - 1 = 171\%$$
            $$EAR = (1+1.71)^{1/10} - 1 = 10.5\%$$
             $$\frac{P_{10}}{P_0} = e^{10r_{cc}}$$
             $$ln(\frac{P_{10}}{P_0}) = 10r_{cc}$$
             $$1 = 1r_{cc}$$
             $$r_{cc} = 10\%$$
    \end{itemize}
    \end{shaded}


\section{Lecture 3: Ch 5, 2: Distributions}

    \begin{shaded}
    \subsection{Practice: Computing Mean and Std Dev}
      \begin{enumerate}
      \setcounter{enumi}{0}
          \item Determine the standard deviation of a random variable $q$ with the following probability distribution:

        \begin{tabular}{lr}
          \toprule
          Value of $q$ & Probability \\
          \midrule
          0 & 0.25 \\
          1 & 0.25 \\
          2 & 0.5 \\
          \bottomrule
        \end{tabular}

        $$\mu_q = 0\cdot0.25 + 1\cdot0.25 + 2\cdot0.5 = 1.25$$
        $$\sigma^2_q = (0-1.25)^2\cdot0.25 + (1-1.25)^2\cdot0.25 + (2-1.25)^2 \cdot 0.5 = 0.6875$$
        $$\sigma = \sqrt{0.6875} = 0.829$$

        \item Your expectations regarding the stock price are as follows:\\
        \vspace{1em}
        
            \begin{tabular}{lrrr}
              \toprule
              Market & Probability & Ending Price & HPR\\
              \midrule
              Boom & 0.35 & \$140 & 44.5\%\\
              Normal & 0.3 & 110 & 14\\
              Recession & 0.35 & 80 & -16.5\\
              \bottomrule
            \end{tabular}
            
        Using your expectations, compute the mean and standard deviation of the HPR for this stock

        $$\mu = 44.5\cdot0.35 + 14\cdot0.3 + -16.5\cdot0.35 = 14\%$$
        $$\sigma^2 = (30.5)^2\cdot0.35 + 30.5^2\cdot0.35 = 930\cdot0.70 = 651$$
        $$\sigma = \sqrt{651} = 25.52\%$$
        
      \end{enumerate}
  \end{shaded}

    \begin{shaded}
    \subsection{Practice: Annualizing Returns}
      \begin{enumerate}
      \setcounter{enumi}{0}
        \item A real-estate property is expected to yield 2\% per quarter with a 
        standard deviation of the quarterly rate of 10\%.  What is the estimated return and standard deviation on the real estate investment after 10 years?  Assuming a normal distribution, what is the probability of loss on the real estate investment after 10 years?

        $$\mu = (1+0.02)^{40} = 2.2 = 120\%$$
        $$\sigma = \sqrt{40}\cdot10 = 6.32\cdot10 = 63.2$$
        $$z=\frac{0-120.8}{63.25}=-1.91$$
        $$P(R<0)=P(Z<-1.91)\approx2.8\%$$
        
        \item The continuously compounded annual return on a stock is normally distributed with a mean of 20\% and standard deviation of 30\%.   With 95.44\% confidence, what pair of values should we expect its actual return in any particular year to be?

        $$20 \pm 2\cdot30 = -30,+80$$
        
      \end{enumerate}
    \end{shaded}

\section{Lecture 4: Ch 5, 3: Evaluation}

    \begin{shaded}
        \subsection{Practice: Fisher Approximation}
        
        \begin{enumerate}
        \item The Fisher equation tells us that the real interest rate approximately
        equals the nominal rate minus the inflation rate.   Suppose the inflation 
        rate increases from 3\% to 5\%.   Does the Fisher equation imply that this 
        increase will result in a fall in the real interest rate?   Explain.\\
         
        \textit{It depends on what happens to nominal interest rates.   If nominal interest rates are unchanged, then yes the real rate must decrease by about 2\%.   However, nominal interest rates are likely to go up a bit, which means real interest rates will fall by less.}
        
        \item Suppose the expected inflation rate rises to 10\%, but the real rate is unchanged.   What happens to the nominal interest rate?\\
        
        \textit{The nominal interest rate must increase.}
        \item The Narnian stock market had a rate of return of 45\% last year, but the inflation rate was 30\%. Using the Fisher Equation, what is the estimated real rate of return for investors.   What was the \textit{actual} real rate of return to Narnian investors?\\
        
        Fisher Approximation: $45\% - 30\% = 15\%$\\
        Actual = $1.45/1.3 = 11\%$
        \end{enumerate}
    \end{shaded}

    \begin{shaded}  
       \subsection{Practice: Excess Returns}

You invest \$27,000 in a corporate bond selling for \$900 per \$1,000 par value.   Over the coming year, the bond will pay interest of \$75 per \$1,000 par value.   The price of the bond at the end of the year will depend on the level of interest rate prevailing at that time.   You construct the following scenario analysis:

        \begin{tabular}{lcc}
            \toprule
             Interest Rates & Probability  & Bond Price \\
             \midrule
             Higher & 0.2 & \$850\\
             Unchanged & 0.5 & 915\\
             Lower & 0.3 & 985\\
             \bottomrule
        \end{tabular}

        Your alternative investment is a T-bill that yields a sure rate of 5\%.   Calculate the HPR for each scenario, the expected rate of return, and the excess rate of return.   What is the expected excess return of the corporate bond?\\
        
        $$\text{Higher: HPR} =  \frac{75-50}{900} = +2.7\% \text{;  excess} = -2.3\%$$
        $$\text{Unchanged: HPR} = \frac{75+15}{900} = 10\% \text{;  excess} =5\%$$
        $$\text{Lower: HPR} = \frac{75+85}{900} = 17.7\% \text{;  excess} =12.7\%$$
        $$\text{Expected Excess} = -2.3\cdot0.2 + 5\cdot0.5 + 12.7\cdot0.3 = 5.85\%$$

    \end{shaded}

\section{Lecture 5: Ch 6, 1: Capital Allocation}

    \begin{shaded}
    \subsection{Practice: Indifference Curves}
    Assume that an investor's utility is
        $$U = E[r] - 0.5 \cdot A \cdot \sigma^2$$ 

    A portfolio has an expected rate of return of 15\% and a standard deviation of 30\%.  T-bills offer a safe rate of return of 2\%.   
    \begin{itemize}
        \item Would an investor with a risk-aversion parameter of $A=4$ prefer to invest in T-bills or the risky portfolio?   \\
        
        \textit{For $A=4$, investing in the risky portfolio gives a utility of $U = 0.15 - 0.5\cdot 4 \cdot \ 0.30^2 = -0.03$, whereas investing in T-Bills gives $U = 0.02$.   T-Bills are preferred.\\}
        \item What if $A=2$?\\
        
        \textit{For $A=2$, investing in the risky portfolio gives a utility of $U = 0.15 - 0.5\cdot 2 \cdot \ 0.30^2 = 0.06$, whereas investing in T-Bills gives $U = 0.02$.   Risky assets are preferred}

        \item Plot the risky portfolio on a graph of $E[r]$ vs. $\sigma$
        \item Assume $A=4$.   How much less return does this investor require to be indifferent between the risky portfolio and an alternative portfolio that has a standard deviation of 25\%?\\

        \textit{5\% less standard deviation, should require less return.   $-2*0.25^2 = -0.125$ v. $-2*0.3^2 = -0.18$.   So need $5.5\%$ less return}
        \item Assume $A=4$.   How much additional risk is this investor willing to bear in order to be indifferent between the risky portfolio and an alternative portfolio that has an expected return of 20\%.\\  

        \textit{First solve for U in the base case.   $U = 0.15 - 0.18 = -0.03$   Now if $E[r]=0.2$ for the same U we can rearrange: $-0.03 = 0.2 - 2*\sigma^2$ then $\sigma^2 = 0.115$ and take a square root, $\sigma=33.9\%$, so an increase of $3.9$ percentage points in standard deviation}
        \item Sketch the indifference curves for $A=4$

        \item How would the indifference curve change for $A=2$\\
    
        \textit{For $A=2$, in the last question, it would take a \textbf{larger} change in the standard deviation for a given change in the expected return.   The indifference curve would be flatter (less curved)}
        
        \item For which value of $A$ is the investor more "risk averse"\\

        \textit{The higher A, the more risk averse}
        
        \item Think about the intuition behind the coefficient $A$ and an investor's risk aversion.
    \end{itemize}

    \end{shaded}

    \begin{shaded}
    \subsection{Practice: Sharpe Ratios}
        \begin{itemize}
        \item Can the Sharpe Ratio of any combination of the risky asset and the risk-free asset be different from the ratio of the risky asset taken alone?\\

        \textit{No.  Rewrite the numerator of the Sharpe Ratio as $w(E[r] - r_f)$, and note that the denominator is $w\sigma$.    The intuition is that scaling the amount invested in the risky asset introduces a $w$ in both the numerator and the denominator of the Sharpe Ratio.   Those $w$s cancel. }
        
        \item The Sharpe Ratio is also referred to as the reward-to-volatility ratio.  Holding the risk-free asset reduces both expected excess return and volatility, but does not change the reward-to-volatility ratio.  Why, then, might an investor choose to hold some of their wealth in the risk-free asset?\\

        \textit{Utility may be higher, even if the Sharpe Ratio is not.   The Sharpe Ratio is the trade off between risk and volatility.   Every investor may value risk and volatility differently.}
    \end{itemize}

    \end{shaded}

\section{Lecture 6: Ch 6, 2: Capital Allocation, Cont'd}   

    \begin{shaded}
    \subsection{Practice: Capital Allocation Problem}
    You manage a risky portfolio with expected rate of return of 12\% and a standard deviation of 28\%.  The T-Bill rate is 2\%.
    \begin{enumerate}
        \item Your client chooses to invest 70\% of a portfolio in your fund and 30\% in a risk-free money 
        market fund.   What are the expected return and standard deviation on her portfolio?
        \begin{align*}
            E[r] &= 0.7 \cdot 0.12 + 0.3 \cdot 0.02 = 0.09 \\
            \sigma &= 0.7 \cdot 0.28  = 0.196
        \end{align*}
        \item What is the reward-to-volatility (Sharpe) ratio of her portfolio?   Of your risky portfolio?   (Hint: draw the CAL, what is the slope?  Where is your client's portfolio?)\\
        $$\text{Sharpe Ratio} = \frac{E[r] - r_f}{\sigma} = \frac{0.07}{0.196} = 0.357$$
        \item Suppose your client instead chooses to invest a proportion $y$ that maximizes the expected return, 
        subject to the constraint that the complete portfolio standard deviation will not exceed 12\%.   What is in the invested proportion $y$?   What is the expected return?
        \begin{align*}
            0.12 &= w*0.28 \implies w=0.428 \\
            E[r] &= 0.428 \cdot 12 + (1-0.428) \cdot 2 = 5.136 + 1.144 = 6.306
        \end{align*}
        \item Your client's degree of risk aversion is $A = 3.5$.   What proportion $y$ should they invest in your fund?

        \begin{align*}
            U &= E[r] - 0.5 \cdot 3.5 \cdot \sigma^2 \\
              &= wE[r_{risky}] + (1-w)r_f - 0.5 \cdot 3.5 w^2\cdot \sigma^2 \\
            &\text{First Derivative wrt w, set equal to 0} \\
             0 &= E[r_{risky}] - r_f - 3.5 \cdot w \cdot \sigma^2 \\
            &\text{Fill in values}\\
             0 &= 0.12 - 0.02 - 3.5 \cdot 0.28^2 \cdot w \\
             w &= \frac{0.10}{3.5*0.28^2}\\
             &= 0.364
        \end{align*}
        
    \end{enumerate}
    \end{shaded}

    \begin{shaded}
    Consider a risky portfolio with $\mathbb{E}[r_p]=8\%$ and $\sigma_p=15\%$.   The risk free rate is 2\%.
    \begin{enumerate}
        \item Your client wants to invest a proportion of her budget in this fund to get a rate of return on her overall portfolio of 5\%.   What proportion of her total portfolio should she invest in the risky fund?
        $$0.05 = w\cdot0.08 + (1-w)\cdot0.02 \implies 0.03 = w \cdot 0.06 \implies w = 0.5$$
        \item What will the standard deviation be on her total portfolio?
        $$7.5\%$$
        \item Another client wants the highest return possible, subject to the constraint that you 
        limit the standard deviation to be no more than 12\%.   Assuming the clients have the same utility functions, except for the parameter $A$, which client is more risk averse?\\
        \textit{The first client is more risk averse (optimal portfolio has less risk)}
    \end{enumerate}
    \end{shaded}

    \begin{shaded}
        \subsection{Practice: Margin Calls}
            Suppose an investor initially pays \$6,000 toward the purchase of \$10,000 worth of stock.   They borrowed the remaining \$4,000 from a broker.  
    \begin{enumerate}
        \item How much equity does the investor begin with in their margin account?   The margin is defined as the $\frac{\text{Equity}}{\text{Value of Stock}}$.   How much margin do they have?\\
            $$\$6,000, 60\%$$
        \item If the stock falls by 30\%, how much margin does the investor have now?\\
            $$\frac{\text{Equity in Account}}{\text{Value of Stock}} = \frac{\$7,000 - \$4,000}{\$7,000} = 42\%$$
        \item If the investor has a maintenance margin of 30\%, what is the maximum amount the stock can fall before they need to post additional margin?   What if there was 50\% maintenance margin?
            \begin{align*}
                &\frac{\$10,000 * (1+r) - \$4,000}{\$10,000 * (1+r)} = 30\% \\
                & -\$4,000 = -\$7,000 \cdot (1+r) \\
                & \implies (1+r) = 0.57 \implies r = -43\%\\
                &\text{For a 50\% maintenance margin}\\
                &\frac{\$10,000 * (1+r) - \$4,000}{\$10,000 * (1+r)} = 50\% \\
                & -\$4,000 = -\$5,000 \cdot (1+r) \\
                & \implies (1+r) = 0.8 \implies r = -20\%\\
            \end{align*}  
        \item What if the investor posts additional margin of 20\% and there is a 10\% maintenance margin?   What \% change in the stock will lead to a margin call?
            \begin{align*}
                &\frac{\$100,000(1+r) - \$80,000}{\$100,000(1+r)} > 0.1\\
                &\$100,000(1+r) - \$80,000 > \$10,000(1+r) \\
                &\$90,000(1+r) > \$80,000\\
                & (1+r) > 8/9\\
                & 1+r > 0.88 \implies r = -12\%
            \end{align*}
    \end{enumerate}

    \end{shaded}

\section{Lecture 7: Ch 7, 1: Diversification}

    \subsection{$E[r]$ and $\sigma$ of a Two Asset Portfolio}

    Construct a portfolio with assets $X$ and $Y$ with returns $r_x$ and $r_y$.  The portfolio is constructed with weights $w_x$ and $w_y$, such that $w_x + w_y = 1$.\\
     
    \textbf{Expected Return: $E[r_p]$}
    \begin{align*}
        E[r_p] &= E[w_xr_x + w_yr_y]\\
                &= E[w_xr_x] + E[w_yr_y]\\
                &= w_xE[r_x] + w_yE[r_y]
    \end{align*}

    \textbf{Standard deviation: $\sigma_p$}\\
    Note that the Variance of X is $Var(X) = E[(X-E[X])^2]$ and $Cov(X,Y) = E[(X-E[X])(Y-E[Y])]$
        \begin{align*}
           \sigma_p^2 = E[(r_p - E[r_p])^2] &= E[(w_xr_x + w_yr_y - E[w_xr_x + w_yr_y])^2]\\
                    &= E[(w_xr_x - E[w_xr_x] + w_yr_y - E[w_yr_y])^2]\\
                    &= E[(w_x(r_x - E[r_x]) + w_y(r_y - E[r_y]))^2]\\
            \text{Let } \tilde{r_x} = r_x - E[r_x] \text{ be the de-meaned return of x}\\
                    & = E[(w_x\tilde{r_x} + w_y\tilde{r_y})^2]\\
                    & = E[w_x^2\tilde{r_x}^2 + 2w_xw_y\tilde{r_x}\tilde{r_y} + w_y^2\tilde{r_y}^2] \\
                \sigma_p^2    & = w_x^2 \operatorname{Var}(r_x) + 2w_xw_y\operatorname{Cov}(r_x,r_y) + w_y^2Var(r_y) \\
                \sigma_p &= \sqrt{w_x^2 \operatorname{Var} + 2w_xw_y\operatorname{Cov}(r_x,r_y) + w_y^2\operatorname{Var}(r_y)}\\
                \text{The correlation coefficient is defined as: }\rho_{x,y} = \frac{\operatorname{Cov}(X,Y)}{\sigma_x\sigma_y}\\
                \sigma_p &= \sqrt{w_x^2 \operatorname{Var} + 2w_xw_y\sigma_x\sigma_y\rho_{x,y} + w_y^2\operatorname{Var}(r_y)}\\
        \end{align*}


    \begin{shaded}
        \begin{enumerate}
        \item The correlation coefficients between pairs of stocks are as follows: Corr(A,B) = 0.85,
        Corr(A,C) = 0.6, Corr(A,D) = 0.45.   Each stock has an expected return of 8\% and a standard 
        deviation of 20\%. 
        \begin{itemize}
            \item If your entire portfolio is now composed of stock A and you can add one stock to your portfolio, which do you add?\\
            
            \textit{Stock D, lower correlation coefficient means more benefit of diversification}
            \item If you compose of portfolio of 0.5 stock A and 0.5 of the stock you chose above, what is the standard deviation of that portfolio?
            $$\sqrt{0.25\sigma_x^2 + 0.25\sigma_y^2 + 2\cdot0.25\sigma_x\sigma_y\rho_{x,y}$$
            $$\sqrt{0.25\cdot0.04 + 0.25\cdot0.04 + 2\cdot0.25\cdot.04\cdot0.45$$
            $$ = \sqrt{0.02+0.19} = 0.17$$
            \textit{Answer: 17\%}
            \item How much did adding this additional stock lower the standard deviation on your portfolio?\\

            \textit{From 20\% to 19.2\%.   Not much}            
            \item Would the answer change for more risk averse investors?\\

            \textit{No, this is not about risk aversion}
            \item Suppose that you could also add T-Bills, which have a rate of 8\%.   How would your portfolio look now?\\

            \textit{All in T-Bills}
        \end{itemize}
    \end{enumerate}
    \end{shaded}

    \begin{shaded}
    \begin{enumerate}
        \item True or false: The standard deviation of the portfolio is always equal to the weighted average of the standard deviations of the assets in the portfolio.\\

        \textit{False, need the covariance term}
    \end{enumerate}
    \end{shaded}


    \begin{shaded}

        \begin{enumerate}
            \item Abigail Grace begins with \$900,000 invested in her original
            portfolio and then inherits \$100,000 of ABC stock. She has the
            following estimates:
            
            \begin{center}
            \begin{tabular}{lrr}
            \toprule
            & $E[r]$, monthly & $\sigma$ \\
            \midrule
            Original Portfolio & 0.67\% & 2.37\% \\
            ABC                 & 1.25\% & 2.95\% \\
            \bottomrule
            \end{tabular}
            \end{center}
            
            The correlation coefficient between ABC and her original portfolio is
            0.40. Answer the following questions:
            
            \begin{itemize}
            
                \item Calculate the expected return of the new portfolio, the
                covariance of ABC's returns with the original portfolio, and the
                standard deviation of the new portfolio.
            
                \textit{The portfolio weights are 0.90 in the original portfolio
                and 0.10 in ABC. The expected return is}
                \[
                    E[r_p]
                    =0.90(0.67\%)+0.10(1.25\%)
                    =0.728\%.
                \]
            
                \textit{The covariance between ABC and the original portfolio is}
                \[
                    \operatorname{Cov}(O,ABC)
                    =\rho_{O,ABC}\sigma_O\sigma_{ABC}
                    =(0.40)(0.0237)(0.0295)
                    =0.00027966.
                \]
            
                \textit{The variance and standard deviation of the new portfolio are}
                \[
                \begin{aligned}
                    \sigma_p^2
                    &=(0.90)^2(0.0237)^2
                    +(0.10)^2(0.0295)^2 \\
                    &\quad
                    +2(0.90)(0.10)(0.00027966)\\
                    &=0.00051401,
                \end{aligned}
                \]
                \[
                    \sigma_p=\sqrt{0.00051401}
                    \approx 0.02267.
                \]
            
                \item Grace sells ABC and invests the proceeds in T-bills that earn
                0.42\% per month. Recalculate the expected return and standard
                deviation of her portfolio.
            
                \textit{Because T-bills are risk-free, their standard deviation and
                covariance with the original portfolio are zero. Therefore,}
                \[
                    E[r_p]
                    =0.90(0.67\%)+0.10(0.42\%)
                    =0.645\%,
                \]
                \[
                    \sigma_p
                    =0.90(2.37\%)
                    =2.133\%.
                \]
            
                \textit{Answer: The portfolio has an expected monthly return of
                0.645\% and a monthly standard deviation of approximately 2.13\%.}
            
                \item Grace is considering selling ABC and buying XYZ, which has the
                same expected return and standard deviation as ABC. Does it matter
                whether Grace owns ABC or XYZ? Why or why not?
            
                \textit{Yes. Expected return and standard deviation alone are not
                sufficient to determine how a security affects a portfolio. Grace
                must also know XYZ's correlation, or covariance, with her original
                portfolio. If XYZ has a lower correlation with the original
                portfolio than ABC does, XYZ will provide a greater diversification
                benefit and produce a lower portfolio standard deviation. If their
                correlations are identical, the two choices will have the same
                portfolio-level expected return and standard deviation.}
            
            \end{itemize}
    \end{enumerate}

    \end{shaded}

        \begin{shaded}
        Consider the following data on Stocks A, B, and C:
    
            \begin{tabular}{lcccc}
            \toprule
            Scenario & Probability & Stock A & Stock B & Stock C\\
            \midrule
            Low & 25\% & 5 & 0 & 12 \\
            High &75\% & 25 & 40 & 24 \\
            \bottomrule
            \end{tabular}
    
        \begin{enumerate}
            \item	What is the expected rate of return for each stock?

            $$E[A] = 0.25\cdot5 + 0.75\cdot25$$, etc
            \item	What is the Variance-Covariance Matrix for Stock A, Stock B, and Stock C (hint: 3x3 matrix)

            $$Cov(A,B) = 0.25\cdot(5-E[A])(0-E[B]) + 0.75\cdot(25-E[A])(40-E[B])$$
            
            \item What is the variance of a portfolio that is composed of 1/3 of each of these stocks?
        \end{enumerate}
    \end{shaded}

\end{document}